\section{A simultaneous bound for lattice intersections}\label{sec:lattice}

A lattice of large index can still contain many vertices of the cube.
The constraint that it contain the columns of a random matrix changes
this picture.  We now show that, with overwhelmingly high probability,
every such lattice has small intersection with every translate of the
cube.  The simultaneous quantifier is essential: the lattices used later
will themselves depend on the matrix.

\begin{theorem}[Simultaneous lattice intersections]\label{thm:small-intersection}
Fix $D>0$ and $0<\alpha<1$, and put $\tau=\ell^{-D}$.
Typically, a nonsingular bit matrix $B$ of size $m\in\mathcal I_N$ has the
following property.  For every lattice $K$ with
\[
 B\Z^m\subseteq K\subseteq\Z^m,
 \qquad [\Z^m:K]\ge N^{\tau N},
\]
and every $x\in\R^m$,
\[
 |(x+K)\cap\bits^m|<2^{\alpha N}.
\]
The exceptional probability is uniform over the bit models specified
in Section~\ref{sec:reductions}.
\end{theorem}

The proof constructs a small collection of vectors $y$ for which $By$
is integral.  A large cube intersection forces these vectors to lie
near a subspace of dimension strictly smaller than $N$.  Their possible
values can therefore be covered by a sufficiently small fixed net.
For fixed representatives, Fourier estimates show that their scalar
products with a random bit row are unlikely to lie simultaneously near
integers.  We test this on disjoint groups of entries in different rows
of $B$; comparing the resulting probability with the net size proves
the theorem without enumerating the lattices $K$.

We first work deterministically.  Assume that $B$ is nonsingular and
satisfies Lemma~\ref{lem:prime-corank}, and suppose that an admissible
$K$ has a coset containing at least $2^{\alpha N}$ cube points.
All constants in the construction are fixed before $N$ tends to
infinity.

\subsection{From a large intersection to separated quotient sums}

Apply Lemma~\ref{lem:minimizer} to $K$ with $T=(\log N)^4$.
Write $L_{\min}$ for the resulting primitive sublattice, and set
\[
 F=\Span L_{\min},\qquad E=F^\perp,\qquad
 \Lambda=\pi_E K,\qquad \Pi=\pi_E\Z^m,
\]
where $\pi_E$ denotes orthogonal projection.  Both projected groups are
lattices in $E$, and $\Lambda\subseteq\Pi$ has full rank there.
Orthogonal factorization of covolumes gives
\begin{equation}\label{eq:lattice-height-index}
 H(F)\le\det L_{\min}\le e^{4N\ell},
 \qquad
 \det\Lambda\ge N^{\tau N}e^{-4N\ell}.
\end{equation}
For the first inequality, $L_{\min}$ is a full-rank sublattice of
$F\cap\Z^m$.  The second bound uses
$\det K=\det L_{\min}\det\Lambda$.

Differences of cube points in the chosen coset belong to $K$ and have
norm at most $\sqrt N$.  Lemma~\ref{lem:projected-count} therefore
bounds the number of their distinct projections into $\Lambda$ by
$2^{o(N)}$.  Some fiber contains $2^{\alpha N-o(N)}$ cube points.
Every affine translate of $F$ contains at most $2^{\dim F}$ cube
points: project injectively onto a set of $\dim F$ independent
coordinate functionals.  Thus
\begin{equation}\label{eq:lattice-large-fiber}
 \dim F\ge\alpha N/2
\end{equation}
for sufficiently large $N$.  This is the dimension saving that will
make the eventual net smaller than the collection of row constraints.

Choose a fixed $P>D+5$.  Since $H(F)\le e^{N\ell^2}$ for large $N$,
Theorem~\ref{thm:height-cover}, with height exponent $2$, gives a
list member $V$ with $d(F,V)\le N\ell^{-P}$.  Put $U=V^\perp$.
Orthocomplementation preserves rank distance, so
$d(E,U)\le N\ell^{-P}$, and $U$ belongs to a fixed list of at most
$2^{N^2\ell^{-P}}$ subspaces.  In particular,
\begin{equation}\label{eq:lattice-cover-dimension}
 \dim U\le(1-\alpha/3)N,
 \qquad W_0=E\cap U,
 \qquad r=\dim E-\dim W_0\le N\ell^{-P}.
\end{equation}
Rank distance controls the dimension of an exact intersection; it
does not assert that the two subspaces have small Euclidean angle.
To obtain vectors close to $U$, we will suppress the remaining $r$
directions by a linear change of scale.

Fix $0<\gamma<\alpha/20$, and choose $C>10$ so large that
\begin{equation}\label{eq:lattice-entropy-margin}
 \Delta:=(C-3)(1-2\gamma)-(C+4)(1-\alpha/3)>0.
\end{equation}
This is possible because $\alpha/3-2\gamma>0$.
Define a positive self-adjoint map $S:E\to E$ by
\begin{equation}\label{eq:lattice-scale}
 S=N^2 I\quad\hbox{on }W_0,
 \qquad
 S=N^{-C-5}I\quad\hbox{on }E\cap W_0^\perp.
\end{equation}
The large scale in $W_0$ will separate quotient sums.  The small scale
on its complement will later force the dual vectors close to $U$.

We need a covolume bound that applies even to sublattices of deficient
rank.  Since $F\cap\Z^m$ is primitive in $\Z^m$,
\begin{equation}\label{eq:lattice-projection-covolume}
 \det\Pi=H(F)^{-1}.
\end{equation}
More generally, every sublattice $L\subseteq\Z^j\oplus\Pi$, for any
$j\ge0$, satisfies
\begin{equation}\label{eq:lattice-lift-bound}
 \det L\ge H(F)^{-1}.
\end{equation}
To see this, take the full inverse image of $L$ in
$\Z^j\oplus\Z^m$ under $(v,z)\mapsto(v,\pi_Ez)$.  Its primitive
kernel is $\{0\}\oplus(F\cap\Z^m)$.  Orthogonal quotient factorization
gives covolume $H(F)\det L$ for the inverse image, which is at least
one because it is an integer lattice.  This argument uses the actual
rank of $L$ and includes rank zero.

The finite group $\Pi/\Lambda$ is a quotient of
$G_B=\Z^m/B\Z^m$, and hence has at most
$g_N=\lceil N^{3/5}\rceil$ generators.  Its order is at least the
lower bound for $\det\Lambda$ in
Equation~\eqref{eq:lattice-height-index}, by
Equation~\eqref{eq:lattice-projection-covolume} and $H(F)\ge1$.
Multiplication by a nonzero integer $u$ has a kernel of size at most
$|u|^{g_N}$ on any group with $g_N$ generators.  Consequently, uniformly
for $1\le |u|\le N^C$,
\begin{equation}\label{eq:lattice-multiple-index}
 |u(\Pi/\Lambda)|\ge N^{\tau N-o(\tau N)}.
\end{equation}
Here $N\ell/\log N=o(\tau N)$ and $N^{3/5}=o(\tau N)$.

For a lattice $L$ in a Euclidean space, write
$\operatorname{dist}(z,L)=\inf_{v\in L}\lVert z-v\rVert_2$.
The next lemma turns the large finite quotient in
Equation~\eqref{eq:lattice-multiple-index} into many separated subset
sums.  Fix an absolute constant $R_0\ge100$.

\begin{lemma}\label{lem:separated-subsets}
For each integer $u$ with $1\le |u|\le N^C$, there is a set
$I_u\subseteq[m]$, of size at least $c_1\tau N$, such that the subset
sums of
\[
 \{uS\pi_E e_i:i\in I_u\}
\]
have pairwise distance greater than $R_0\sqrt N$ modulo $S\Lambda$.
The constant $c_1>0$ is fixed, independently of $u$ and $N$.
\end{lemma}

\begin{proof}
Choose a maximal such set of indices, say $i_1,\ldots,i_j$.
Adding an unselected index creates a pair of subset sums at distance
at most $R_0\sqrt N$ modulo $S\Lambda$.  The added index occurs in
exactly one of them, since the old subset sums were separated.
Undoing $S$, we obtain, for every $i\in[m]$,
\begin{equation}\label{eq:lattice-representations}
 u\pi_Ee_i\equiv
 \sum_{h=1}^j a_{hi}u\pi_Ee_{i_h}+\varepsilon_i
       \pmod\Lambda,
 \qquad a_{hi}\in\{-1,0,1\},
 \qquad\lVert S\varepsilon_i\rVert_2\le R_0\sqrt N,
\end{equation}
where $\varepsilon_i\in\Pi$.  For the selected indices use the exact
representations, with coefficient columns the standard basis and
zero errors.

Let $t$ be the dimension spanned by the errors, and choose $t$
independent ones.  By Equation~\eqref{eq:lattice-lift-bound}, their
volume is at least $1/H(F)$.  On every $t$-dimensional subspace, $S$
expands volume by at least $N^{2t-(C+7)r}$: at most $r$ singular
directions have the smaller scale in Equation~\eqref{eq:lattice-scale}.
Hadamard's inequality for the scaled errors gives
\[
 H(F)^{-1}N^{2t-(C+7)r}\le(R_0\sqrt N)^t.
\]
Using Equation~\eqref{eq:lattice-height-index}, we conclude that
\begin{equation}\label{eq:lattice-error-rank}
 t=O_C\left(r+\frac{N\ell}{\log N}\right)=o(\tau N).
\end{equation}

Define a real linear map $M:\R^m\to\R^j\oplus E$ whose $i$th
column is $(a_{1i},\ldots,a_{ji},\varepsilon_i)$.
Its integer image $M\Z^m$ is a sublattice of $\Z^j\oplus\Pi$,
of rank $s\le j+t$.  Evaluation of a representation defines the
homomorphism
\[
 \Psi:\Z^j\oplus\Pi\longrightarrow\Pi/\Lambda,
 \qquad
 \Psi(a,v)=\sum_{h=1}^j a_hu\pi_Ee_{i_h}+v\pmod\Lambda.
\]
By Equation~\eqref{eq:lattice-representations},
$\Psi(M\Z^m)=u(\Pi/\Lambda)$.
Every column of $MB$ lies in $\ker\Psi$, because its evaluation is
$u\pi_EBe_i\in\Lambda$.  Moreover, $B$ is nonsingular, so
$\rank(MB)=s$.  Choose $s$ independent columns of $MB$, and let
$L'$ be their integer span.  Then
\[
 |u(\Pi/\Lambda)|\le[M\Z^m:L']
       =\frac{\det L'}{\det(M\Z^m)}.
\]
Each error has ordinary norm at most $R_0N^{C+5}\sqrt N$.
Since the entries of $B$ are bounded by one, the chosen columns have
norm at most $N^{C+9}$ for large $N$.  Equation~\eqref{eq:lattice-lift-bound}
and Hadamard's inequality therefore imply
\[
 N^{\tau N-o(\tau N)}
 \le H(F)N^{(C+9)s}
 \le H(F)N^{(C+9)(j+t)}.
\]
Together with Equation~\eqref{eq:lattice-error-rank}, this yields
$j\ge c_1\tau N$, for example with $c_1=1/[2(C+9)]$ and sufficiently
large $N$.
\end{proof}

There are only polynomially many integers $u$ under consideration.
Select each coordinate independently with probability $\gamma/2$.
Binomial concentration, applied to $[m]$ and to the sets $I_u$, shows
that some set $J\subseteq[m]$ satisfies
\begin{equation}\label{eq:lattice-small-coordinate-set}
 |J|\le\gamma N,
 \qquad |J\cap I_u|\ge c_2\tau N
       \quad(1\le|u|\le N^C),
\end{equation}
with a fixed $c_2>0$.  Indeed the failure probability for each $I_u$
is $\exp(-\Omega(\tau N))$, which absorbs the polynomial union.

\subsection{Dual vectors and Fourier bounds}

We have obtained separated quotient sums using only the coordinates
in $J$.  We now convert that separation into vectors whose scalar
products with random bit columns are spread around the circle.
The dual lattice is
\[
 \Lambda^*=\{y\in E:\langle y,\Lambda\rangle\subseteq\Z\}.
\]
For every $y\in\Lambda^*$, symmetry of $B$ and
$\pi_E B\Z^m\subseteq\Lambda$ give
\begin{equation}\label{eq:lattice-dual-integrality}
 By\in\Z^m.
\end{equation}
Sample $y\in\Lambda^*$ with probability proportional to
$\exp(-\pi\lVert S^{-1}y\rVert_2^2)$.
The factor $S^{-1}$ strongly suppresses components in
$E\cap W_0^\perp$; Gaussian tails will therefore place the sampled
vectors close to $W_0\subseteq U$.  Write
\[
 \varphi(t)=b+p e^{2\pi i t}.
\]

\begin{lemma}\label{lem:dual-fourier-average}
There is a fixed $c_3>0$ such that, for every integer
$1\le|u|\le N^C$ and every collection of real numbers $(s_i)_{i\in J}$,
\begin{equation}\label{eq:lattice-fourier-average}
 \E_y\prod_{i\in J}|\varphi(uy_i+s_i)|^2
       \le e^{-c_3\tau N}.
\end{equation}
\end{lemma}

\begin{proof}
Let $X=\xi-\xi'$ where $\xi,\xi'$ are independent vectors of
Bernoulli $p$ bits on $J$, extended by zero outside $J$.
Expanding the product in Equation~\eqref{eq:lattice-fourier-average}
expresses its expectation as
\[
 \E_X\left[e^{2\pi i\sum_{i\in J}s_iX_i}
                  \E_y e^{2\pi i u\langle y,X\rangle}\right].
\]
Poisson summation in $E$ gives
\begin{equation}\label{eq:lattice-poisson-ratio}
 \E_y e^{2\pi i u\langle y,X\rangle}=f(uX),\qquad
 f(uX)=
 \frac{\sum_{\lambda\in\Lambda}
       e^{-\pi\lVert S(\lambda-u\pi_EX)\rVert_2^2}}
      {\sum_{\lambda\in\Lambda}e^{-\pi\lVert S\lambda\rVert_2^2}}.
\end{equation}
The ratio is nonnegative and at most one.  Positivity follows from
the displayed formula; the upper bound follows either from its
characteristic-function interpretation or from the Gaussian
translate inequality.  Thus taking absolute values permits us to
discard the phases involving $s_i$ and bound the original expectation
by $\E_X f(uX)$.

We record explicitly the Gaussian estimates used here.  For every
lattice $L$ of rank $d$ and every $z$ in its span, Poisson summation
and positivity of the Fourier transform of a Gaussian imply
\[
 \sum_{v\in L}e^{-\pi\lVert v-z\rVert_2^2/2}
 \le\sum_{v\in L}e^{-\pi\lVert v\rVert_2^2/2},
 \qquad
 \frac{\sum_{v\in L}e^{-\pi\lVert v\rVert_2^2/2}}
      {\sum_{v\in L}e^{-\pi\lVert v\rVert_2^2}}
 \le 2^{d/2}.
\]
For the second inequality, Poisson summation expresses the ratio as
$2^{d/2}$ times a ratio of dual Gaussian masses at decreasing scales.
These estimates hold for arbitrary lattices, with no additional
regularity assumption.  For these lattice Gaussian comparisons, see
Banaszczyk~\cite{Banaszczyk} and
\cite[Equation (5) and Claim 2.2]{ReverseMinkowski}.

Apply them to $L=S\Lambda$.  If
$\operatorname{dist}(uS\pi_EX,S\Lambda)\ge R_0\sqrt N/3$, retaining
half of each Gaussian exponent in the numerator of
Equation~\eqref{eq:lattice-poisson-ratio} gives
\begin{equation}\label{eq:lattice-gaussian-distance}
 f(uX)\le e^{-\pi R_0^2N/18}2^{\dim E/2}
          \le e^{-c_4N}
\end{equation}
for an absolute $c_4>0$.

To bound the complementary event, condition on $\xi'$ and on all
coordinates of $\xi$ except those in $J\cap I_u$.
By Lemma~\ref{lem:separated-subsets}, at most one of the remaining
subset sums lies within distance $R_0\sqrt N/3$ of the prescribed
translate of $S\Lambda$.  Otherwise the distance between two such
sums would be at most $2R_0\sqrt N/3$.
The probability of any one bit pattern is at most
$b^{|J\cap I_u|}\le b^{c_2\tau N}$.  Combining this bound with
Equation~\eqref{eq:lattice-gaussian-distance} proves the result,
after decreasing $c_3$.
\end{proof}

We next choose several dual vectors at once.  Let
\begin{equation}\label{eq:lattice-number-testers}
 k=\left\lfloor\kappa\frac{\tau N}{\log N}\right\rfloor,
\end{equation}
where $\kappa>0$ will be fixed sufficiently small.
For fixed $c_5>0$ small enough, there exist
$y^1,\ldots,y^k\in\Lambda^*$ such that
\begin{equation}\label{eq:lattice-tester-geometry}
 \lVert y^g\rVert_2\le R_0N^2\sqrt N,
 \qquad\operatorname{dist}(y^g,U)\le R_0N^{-C-5}\sqrt N
 \qquad(1\le g\le k),
\end{equation}
\begin{equation}\label{eq:lattice-tester-fourier}
 \prod_{i\in J}\left|\varphi\left(\sum_{g=1}^k w_gy_i^g\right)\right|
 \le e^{-c_5\tau N}
 \quad\bigl(w\in\Z^k,\ 1\le\lVert w\rVert_\infty\le N^C\bigr).
\end{equation}
Here and below a distance to a subspace is Euclidean distance.

For completeness, sample the $k$ vectors independently according to
the dual Gaussian.  The Gaussian scale inequality just proved,
applied on $S^{-1}\Lambda^*$, gives
\[
 \Prob\bigl(\lVert S^{-1}y^g\rVert_2>R_0\sqrt N\bigr)
 \le 2^{\dim E/2}e^{-\pi R_0^2N/2}.
\]
Outside this event, the two scales of $S$ give
Equation~\eqref{eq:lattice-tester-geometry}; for the second bound use
$W_0\subseteq U$.  For a fixed nonzero $w$ in the stated range,
condition on every vector except one whose coefficient is nonzero.
The remaining vectors produce arbitrary shifts in
Lemma~\ref{lem:dual-fourier-average}.  Markov's inequality bounds
failure of Equation~\eqref{eq:lattice-tester-fourier} by
\[
 \exp\bigl(-(c_3-2c_5)\tau N\bigr).
\]
There are at most $(3N^C)^k$ possible $w$.
Choose $c_5<c_3/4$ and then $\kappa$ so small that
\begin{equation}\label{eq:lattice-kappa-choice}
 (C+2)\kappa<c_5/2.
\end{equation}
The union of these failures and the norm failures has probability
less than one for large $N$, establishing the asserted existence.
Notice the order of choices: $\gamma,C$ were fixed first, the
constants $c_1,c_2,c_3,c_5$ follow, and only then is $\kappa$ chosen.
In particular, the dependence of the Fourier constants on $C$ does
not affect the positive margin in
Equation~\eqref{eq:lattice-entropy-margin}.

We have now associated to every violating lattice a set $J$ of at
most $\gamma N$ coordinates and $k$ vectors satisfying
Equations~\eqref{eq:lattice-tester-geometry}--\eqref{eq:lattice-tester-fourier},
together with $By^g\in\Z^m$.  The first two conditions will allow a
fixed finite description; the last condition will make such a
description unlikely for a random matrix.

\subsection{Counting fixed descriptions and testing rows}

Fix a covering-list member $U$ with
$\dim U\le(1-\alpha/3)N$.  The region for one vector in
Equation~\eqref{eq:lattice-tester-geometry} can be covered by at most
\begin{equation}\label{eq:lattice-net-size}
 N^{(C+4)\dim U}
\end{equation}
sets of diameter $4N^{-C}$.  Indeed cover its projection to $U$, a
ball of radius at most $R_0N^{5/2}$, at scale $N^{-C}$.
The entire perpendicular thickness is
$2R_0N^{-C-9/2}=o(N^{-C})$, so it requires no additional covering
directions.  This also proves the assertion when $\dim U=0$, with
one set.

Fix $J$ and one such cell for each of the $k$ vectors.  If the cells
admit a tuple satisfying
Equations~\eqref{eq:lattice-tester-geometry}--\eqref{eq:lattice-tester-fourier},
choose one such tuple and denote it by $\widehat y^1,\ldots,
\widehat y^k$.  These choices are made independently of $B$: the
conditions just imposed involve only the fixed $U,J$ and the vector
coordinates.  The representatives need not belong to any dual
lattice.  Nevertheless, existence of an actual witness in the same
cells implies, for every row index $v$ and every $g$,
\begin{equation}\label{eq:lattice-approximate-congruences}
 \operatorname{dist}\bigl((B\widehat y^g)_v,\Z\bigr)
       \le 4N^{-C}\sqrt N,
\end{equation}
because $By^g$ is integral and $\lVert B_{v,*}\rVert_2\le\sqrt N$.
This passage is what removes the dependence on the unknown lattice
from the probability estimate.

Condition on every entry of $B$ except those in the cross block
between $J$ and $J^c$.  For a row $v\in J^c$, the remaining variables
are independent bits $(B_{vi})_{i\in J}$.  For different $v\in J^c$
these are disjoint collections of upper-triangular variables, so the
row events in Equation~\eqref{eq:lattice-approximate-congruences}
are conditionally independent.

We bound the probability of one row event by a finite Fourier
expansion.  Put $M_0=\lfloor N^{C-2}\rfloor$, and define the
nonnegative periodic kernel
\[
 F_{M_0}(t)=\left|\frac1{M_0}
                   \sum_{h=0}^{M_0-1}e^{2\pi i ht}\right|^2.
\]
On the interval of radius $4N^{-C}\sqrt N$ around an integer this
kernel is bounded below by a fixed positive constant, since the
product of that radius and $M_0$ tends to zero.
Its constant Fourier coefficient is $M_0^{-1}$, its other
coefficients have magnitude at most $M_0^{-1}$, and its frequencies
satisfy $|w|<M_0$.  Apply the product of $k$ such kernels to the row
vector $((B\widehat y^g)_v)_{g=1}^k$.

For every nonzero frequency vector $w$, the absolute value of its
conditional expectation is bounded by
\[
 \prod_{i\in J}
 \left|\varphi\left(\sum_{g=1}^k w_g\widehat y_i^g\right)\right|
 \le e^{-c_5\tau N}.
\]
The deterministic contribution from $J^c$ has modulus one.  A bit
with success probability $b$ has the same characteristic-function
modulus as one with success probability $p$, since
$|p+b e^{it}|^2=|b+p e^{it}|^2$.
All frequencies in the product kernel fall within
Equation~\eqref{eq:lattice-tester-fourier}.
It follows, for an absolute constant $C_2$, that the row probability
is at most
\begin{align}
 C_2^kM_0^{-k}\left(1+(3M_0)^k e^{-c_5\tau N}\right)
 &\le N^{-(C-3)k}.
 \label{eq:lattice-row-cost}
\end{align}
The last inequality uses Equation~\eqref{eq:lattice-kappa-choice};
the quantity inside parentheses is bounded, and fixed factors per
coordinate are absorbed by one power of $N$.

There are $m-|J|\ge(1-2\gamma)N$ tested rows for large $N$.
The conditional independence just noted and
Equation~\eqref{eq:lattice-row-cost} therefore bound the probability
for a fixed $U,J$, and list of cells by
$N^{-(C-3)(1-2\gamma)kN}$.
The cell choices cost at most
$N^{(C+4)(1-\alpha/3)kN}$ by Equation~\eqref{eq:lattice-net-size}.
There are at most $2^m$ choices of $J$, and at most
$2^{N^2\ell^{-P}}$ choices of $U$.
Consequently the probability of any witness, together with
nonsingularity and the prime-corank bound, is at most
\[
 2^{m+N^2\ell^{-P}}
       \exp(-\Delta kN\log N)
 \le \exp(-c_6\tau N^2)
\]
for a fixed $c_6>0$.  Here $k\sim\kappa\tau N/\log N$,
$\Delta>0$, and $P>D$.
Adding the exceptional probability from
Lemma~\ref{lem:prime-corank} still gives $2^{-\omega(N)}$.
Every violating lattice supplied a witness in the preceding
construction, so this proves Theorem~\ref{thm:small-intersection}.
